\documentclass{article}

% Language setting
% Replace `english' with e.g. `spanish' to change the document language
\usepackage[english]{babel}

% Set page size and margins
% Replace `letterpaper' with `a4paper' for UK/EU standard size
\usepackage[letterpaper,top=2cm,bottom=2cm,left=3cm,right=3cm,marginparwidth=1.75cm]{geometry}
\usepackage{natbib}
\usepackage{booktabs}


% ============================================================
% Figures generated by the analysis pipeline
%
% main.tex lives in latex/, while generated figures live in
% figures/current/.
%
% figures/current is a stable symlink to the current
% provenance-specific draft directory.
% ============================================================
\usepackage{graphicx}

% ============================================================
% FIGURE PATHS
%
% Generated figures are searched before legacy snapshots so
% that re-running the analysis automatically updates the
% manuscript at the next LaTeX compilation.
%
% Compile from the repository root with:
%
%     latexmk -pdf -cd latex/main.tex
%
% ============================================================
\graphicspath{
    {../figures/current/}
    {../results/paper_figures/}
    {../results/bias_paper_figures/}
    {../figures/finite_source_sanity/}
    {../figures/roman_asimov/}
    {../figures/roman_asimov/representative_lightcurves/}
    {../figures/appendix/}
    {figures/}
}


% Useful packages
\usepackage{amsmath}
\usepackage{float}
\usepackage{pdflscape}
\usepackage{amssymb}

\usepackage[colorlinks=true, allcolors=blue]{hyperref}
\usepackage{subcaption}

\title{Intrinsic Degeneracies Between Binary-Source and Point-Source Microlensing Events}

\author{A. Varela, R. Di Stefano, S. Kahkpash, M. Makler, C. Galmes Altafulla}

\begin{document}
\maketitle


\begin{abstract}
Binary-source microlensing events produced by a single point lens (BSPL)
can closely resemble standard point-source point-lens (PSPL) events when
source orbital motion is absorbed by shifts in the fitted PSPL
parameters. We quantify this intrinsic degeneracy by fitting simulated
binary-source light curves with re-optimized PSPL models and measuring
the normalized residual mismatch $D_{\rm BSPL-PSPL}$.

For a single luminous source, the mismatch is controlled jointly by the
orbital displacement and by the fraction of the source orbit sampled over
the microlensing event. Small residuals can coexist with biased PSPL
parameters; for the fiducial $t_E=30\,$d calculation, the fitted $u_0$ and
$t_E$ shifts are strongly correlated. When both sources are luminous, the
first-order displacement of the unresolved source photocenter vanishes on
the analytic locus $q_f=q_M$. A Taylor expansion and direct numerical tests
show that the leading small-separation mismatch changes from
$D_{\rm BSPL-PSPL}\propto\xi_{\rm rel}$ for a dark companion to
$D_{\rm BSPL-PSPL}\propto\xi_{\rm rel}^{2}$ on this cancellation locus.
A PARSEC main-sequence calculation in the Roman F146 band shows that
$q_f=q_M$ is not a generic stellar mass--flux relation, but that strong
partial photocenter suppression is physically accessible for nearly
equal-mass binaries.

We finally propagate the matched $t_E=30\,$d single-luminous-source grid
through a Roman-like W149 Asimov experiment. The complete survey-weighted
calculation contains 23,940 PSPL fits for W149$=19$, 21, and 23. The Roman
maps broadly preserve the intrinsic degeneracy structure while shifting
the model-separation boundary with source brightness; a fixed intrinsic
mismatch therefore does not correspond to a universal survey detection
threshold. The reported Roman $\Delta\chi^2$ values are expected model
separations for the adopted cadence and noise prescription, not calibrated
$p$-values or population-level event fractions.
\end{abstract}


\section{Introduction}
\label{sec:introduction}

Gravitational microlensing probes compact objects and exoplanets in
regions of the Galaxy that are difficult to access with other
techniques \citep{MaoPaczynski1991,BennettRhie1996}. The standard
point-source point-lens (PSPL) model assumes a single point source moving
rectilinearly relative to a single point lens, but source multiplicity
and source orbital motion can perturb this simple geometry.

Binary-source microlensing was recognized early as a distinct class of
single-lens events \citep{GriestHu1992}. If both components are
luminous, the observed flux is the flux-weighted sum of the two
independently magnified sources. Early studies emphasized that a large
fraction of such events can remain difficult to identify because of
unequal source fluxes, unfavorable source separations, blending, or
because the resulting light curve can mimic a single-source event with
different fitted parameters
\citep{1998A&A...333..893D,HanJeong1998}. This already makes the
binary-source problem a model-degeneracy problem rather than simply a
question of whether two stars are physically present.

Binary-source signals can also mimic perturbations that would otherwise
be attributed to more complex lens models. In particular, a subset of
binary-source light curves can resemble planetary microlensing
perturbations \citep{Gaudi1998,2002ApJ...564.1015H}. This ambiguity
continues to be relevant in modern high-cadence data. For example,
binary-source and planetary interpretations competed directly in
MOA-2012-BLG-486, while the planetary candidate MOA-2006-BLG-074 was
ultimately better explained by source orbital motion
\citep{Hwang2013,2021AJ....162...59R}. Recent work has also revisited
the binary-source degeneracy explicitly in the context of wide-orbit
planets and survey cadence \citep{2025AJ....170..141S}.

A related but physically distinct manifestation of source binarity is
orbital motion. If one component is effectively dark, the observable
effect is the motion of the luminous source around the source-binary
barycenter; more generally, when both components contribute light, the
relevant first-order displacement is that of the unresolved source
photocenter. Orbital source motion is commonly referred to as
\emph{xallarap}. It has been inferred in observed microlensing events
\citep{Alcock2001} and can imitate annual microlensing parallax
\citep{Poindexter2005}. Source orbital motion has also been explored as
a way to detect unseen companions and planets
\citep{RahvarDominik2009,Miyazaki2021}, including applications to
future space-based microlensing surveys. A complementary binary-source
channel in which a luminous planetary companion contributes directly to
the observed flux has likewise been investigated for WFIRST/Roman-like
surveys \citep{2019MNRAS.490.1581B}.

These previous studies therefore contain two closely related limits that
are often treated separately: the flux-weighted superposition of two
luminous source stars and the orbital displacement of a luminous source
or unresolved source photocenter. In this work we use
\emph{binary-source} for the physical source system and treat both
limits within the same BSPL framework. The lens is taken to be
point-like throughout. We use \emph{xallarap} specifically when
referring to source orbital motion, most directly in the
single-luminous-source or dark-companion limit.

Our focus is also different from asking whether a binary-source or
xallarap model can provide a statistically significant improvement over
PSPL, or whether the orbital parameters can be uniquely recovered.
Instead, we ask a complementary model-space question: after the standard
PSPL parameters are re-optimized, how close can a true BSPL light curve
remain to the PSPL model family?

The analysis makes three specific contributions. First, we treat the
best-fitting PSPL solution as a projection of the BSPL light curve onto
the re-optimized PSPL model family, rather than comparing the BSPL signal
with an unperturbed PSPL trajectory. Second, we separate the
binary-source perturbation into an \emph{absorbed component}, quantified
by shifts in $(t_0,u_0,t_E)$, and an \emph{irreducible component},
quantified by the residual mismatch after PSPL re-optimization. Third,
for two luminous sources we derive the analytic first-order
photocenter-cancellation locus $q_f=q_M$ and show analytically and
numerically that it changes the leading small-separation mismatch from
linear to quadratic order in the relative orbital displacement. We then
separate this geometric limiting case from its astrophysical accessibility
by evaluating Roman-band flux ratios from a representative main-sequence
stellar isochrone.

We quantify the irreducible component with the normalized distance
$D_{\rm BSPL-PSPL}$ and study the corresponding PSPL parameter biases.
We first isolate the intrinsic geometry of the BSPL--PSPL degeneracy and
then propagate the same configurations through an adopted Roman cadence
and W149 uncertainty model using noiseless Asimov data. This separation
allows us to distinguish a physical model-space degeneracy from the
survey-dependent significance with which the residual can be measured.

Section~\ref{sec:simulations} summarizes the binary-source model and the
physical scales used to organize the experiments.
Section~\ref{sec:degeneracy_metrics} defines the intrinsic mismatch and
parameter offsets. Sections~\ref{sec:one_luminous_source} and
\ref{sec:two_luminous_sources} present the single- and two-luminous-source
limits, respectively. Section~\ref{sec:roman_asimov} translates the
single-luminous-source grid into Roman Asimov model separation. The
complete orbital formalism is given in Appendix~\ref{app:bspl_formalism}.
Numerical robustness tests, including sensitivity to the integration
window, are collected in Appendix~\ref{app:numerical_robustness}.
The dependence on projected orbital geometry is tested in
Appendix~\ref{app:geometry_robustness}.


\section{Binary-source simulations}
\label{sec:simulations}

We simulate a point lens and a binary source on a circular Keplerian
orbit. The barycentric microlensing trajectory is described by the
standard parameters $(t_0,u_0,t_E)$. The source masses are
$M_{S,1}$ and $M_{S,2}$, with mass ratio
\begin{equation}
q_M=\frac{M_{S,2}}{M_{S,1}}.
\label{eq:qM}
\end{equation}
For orbital period $P$, Kepler's law fixes the relative semimajor axis,
\begin{equation}
a_{\rm rel}^{3}
=
\frac{G(M_{S,1}+M_{S,2})P^{2}}{4\pi^{2}}.
\label{eq:kepler_source}
\end{equation}
and hence the dimensionless relative orbital amplitude
\begin{equation}
\xi_{\rm rel}
=
\frac{a_{\rm rel}}{\hat r_E},
\qquad
\hat r_E=D_S\theta_E.
\label{eq:xi_rel}
\end{equation}
The orbital amplitude of source 1 about the barycenter is
\begin{equation}
\xi_{E,1}
=
\frac{q_M}{1+q_M}\,\xi_{\rm rel}.
\label{eq:xi_E1}
\end{equation}
In the single-luminous-source limit we denote this quantity by
$\xi_E$.

Two scale comparisons are particularly useful. Spatially,
\begin{equation}
\frac{\xi_E}{u_0}
\end{equation}
compares the orbital displacement with the minimum barycentric
lens--source separation. Temporally, $P/t_E$ compares the orbital period
with the global event timescale, while
\begin{equation}
\frac{P}{t_{\rm eff}}
=
\frac{P}{u_0t_E},
\qquad
t_{\rm eff}=u_0t_E.
\label{eq:teff}
\end{equation}
is useful near the high-magnification region. These ratios motivate the
coordinates used below, but they do not completely determine the light
curve because the projected orbital geometry also matters.

The normalization by $\hat r_E$ connects the same physical source binary
to different microlensing configurations. For fixed lens--source
geometry,
\begin{equation}
\hat r_E\propto M_L^{1/2},
\qquad
\xi_E\propto M_L^{-1/2},
\qquad
t_E\propto M_L^{1/2},
\end{equation}
so changing the lens mass modifies both the spatial and temporal
dimensionless scales.

The projected orbit is specified by inclination $i$, phase $\phi$, and
orientation $\alpha$. These angles are held fixed unless explicitly
varied. When both sources are luminous we define
\begin{equation}
q_f=\frac{F_{S,2}}{F_{S,1}}.
\label{eq:qf}
\end{equation}
with $q_f=0$ corresponding to a single luminous source. The full
trajectories and combined magnification are given in
Appendix~\ref{app:bspl_formalism}.

The simulations therefore map physically motivated binary-source
parameters into magnification light curves. The non-trivial step is the
subsequent projection of each light curve onto the PSPL model family,
which we describe next.

For convenience, Table~\ref{tab:notation} collects the notation used
throughout the analysis.

\begin{table}[t]
\centering
\small
\caption{Principal notation used throughout the paper.}
\label{tab:notation}
\begin{tabular}{p{0.20\linewidth}p{0.57\linewidth}p{0.15\linewidth}}
\toprule
Symbol & Definition & First definition \\
\midrule
$t_0,u_0,t_E$ & PSPL time of closest approach, impact parameter, and Einstein timescale & Sec.~\ref{sec:simulations} \\
$M_{S,1},M_{S,2}$ & Source-component masses & Sec.~\ref{sec:simulations} \\
$q_M$ & Source mass ratio, $M_{S,2}/M_{S,1}$ & Eq.~\ref{eq:qM} \\
$q_f$ & Unmagnified source-flux ratio, $F_{S,2}/F_{S,1}$ & Eq.~\ref{eq:qf} \\
$P$ & Source-binary orbital period & Eq.~\ref{eq:kepler_source} \\
$a_{\rm rel}$ & Relative source-binary semimajor axis & Eq.~\ref{eq:kepler_source} \\
$\hat r_E$ & Einstein radius projected to the source plane & Eq.~\ref{eq:xi_rel} \\
$\xi_{\rm rel}$ & Relative orbital separation in units of $\hat r_E$ & Eq.~\ref{eq:xi_rel} \\
$\xi_{E,1}$, $\xi_E$ & Orbital amplitude of source 1; denoted $\xi_E$ in the single-luminous limit & Eq.~\ref{eq:xi_E1} \\
$t_{\rm eff}$ & High-magnification timescale $u_0t_E$ & Eq.~\ref{eq:teff} \\
$\alpha,\phi,i$ & Sky-plane orientation, orbital phase, and inclination & App.~\ref{app:bspl_formalism} \\
$D_{\rm BSPL-PSPL}$ & Intrinsic normalized mismatch after PSPL re-optimization & Eq.~\ref{eq:D_metric} \\
$\Delta\chi^2_{\rm Roman}$ & Roman Asimov PSPL--BSPL model separation & Eq.~\ref{eq:roman_dchi2} \\
$D_{\rm Roman,w}$ & Survey-weighted mismatch $\sqrt{\Delta\chi^2_{\rm Roman}}/{\rm SNR}_{\rm event}$ & Eq.~\ref{eq:D_roman} \\
\bottomrule
\end{tabular}
\end{table}

\begin{table}[t]
\centering
\caption{Principal numerical setup used for the production scans. We adopt
$t_E=30\,$d as the fiducial timescale for the main intrinsic maps, representative
of ordinary Galactic-bulge microlensing events. The multi-$t_E$ comparison
explicitly retains longer-duration cases, including $t_E=150\,$d, and some
auxiliary validation calculations are left at their original long-timescale
setup where stated.}
\label{tab:numerical_setup}
\begin{tabular}{p{0.36\linewidth}p{0.56\linewidth}}
\toprule
Quantity & Adopted value or range \\
\midrule
$M_{S,1},M_{S,2}$ & $2\,M_\odot,1\,M_\odot$ \\
$q_M$ (fiducial one-luminous-source scan) & $0.5$ \\
$\hat r_E$ & $5\,$AU \\
$t_E$ (fiducial) & $30\,$d \\
$u_0$ grid & $10^{-2}$--$10$, 200 logarithmic nodes \\
$P$ grid & $10$--$10^5\,$d, 60 logarithmic nodes \\
Intrinsic time window & $t_0\pm3.5\,t_E$, $10^4$ uniformly sampled epochs \\
$(\alpha,\phi,i)$ & $(0,0,\pi/2)$ unless stated otherwise \\
One-luminous $q_M$ scan & $10^{-4}$--$1$, 300 logarithmic nodes at fixed $u_0=0.1$ and fixed $M_{S,1}+M_{S,2}=3\,M_\odot$ \\
$P$ grid for one-luminous $q_M$ scan & $10$--$10^5\,$d, 300 logarithmic nodes \\
Two-luminous $q_M$ grid & $10^{-4}$--$1$, 200 logarithmic nodes, with $q_M=M_{S,2}/M_{S,1}$ and fixed $M_{S,1}+M_{S,2}=3\,M_\odot$ \\
Two-luminous $q_f$ grid & $0$ plus 200 logarithmic nodes in $10^{-4}$--$1$ \\
Two-luminous $P/t_E$ & $0.3,1,3$ \\
Intrinsic PSPL fit & free $(t_0,u_0,t_E)$ in magnification space; Nelder--Mead minimization of the trapezoidal intrinsic objective \\
\bottomrule
\end{tabular}
\end{table}

The scans use two complementary mass conventions. In the fiducial
Kepler-consistent family we keep $M_{S,1}=2\,M_\odot$ and
$M_{S,2}=1\,M_\odot$ fixed, so that changing $P$ changes the physical
separation through Kepler's law. In all scans in which $q_M$ is varied
independently, including the one-luminous $q_M$ scan and the
two-luminous $(q_M,q_f)$ grids, we instead keep the total source mass
fixed at $M_{\rm tot}=M_{S,1}+M_{S,2}=3\,M_\odot$. With
$q_M=M_{S,2}/M_{S,1}$, the component masses are therefore
$M_{S,1}=M_{\rm tot}/(1+q_M)$ and
$M_{S,2}=q_M M_{S,1}$. This keeps $a_{\rm rel}(P)$ fixed at a given
period and isolates the effect of redistributing the same total source
mass between the two components.


\section{Intrinsic BSPL--PSPL mismatch}
\label{sec:degeneracy_metrics}

For each simulated BSPL light curve, we fit a PSPL model by optimizing
$(t_0,u_0,t_E)$ in magnification space. The production minimization is
performed with the Nelder--Mead algorithm, initialized at the true
barycentric $(t_0,u_0,t_E)$ values, using the same objective that defines
the numerator of $D_{\rm BSPL-PSPL}$. All intrinsic fits and metrics are
evaluated over the fixed event-centered window
$[t_0-3.5t_E,t_0+3.5t_E]$, sampled with $10^4$ uniformly spaced epochs.
The best-fitting model minimizes
\begin{equation}
\mathcal{J}
=
\int_{t_0-3.5t_E}^{t_0+3.5t_E}
\left[
A_{\rm BSPL}(t)-A_{\rm PSPL}(t)
\right]^2
\,dt,
\end{equation}
with the integral evaluated numerically using trapezoidal integration.
We define the residual
\begin{equation}
\Delta A(t)
=
A_{\rm BSPL}(t)-A_{\rm PSPL,best}(t),
\end{equation}
and use as our primary intrinsic metric
\begin{equation}
D_{\rm BSPL-PSPL}
=
\left[
\frac{
\displaystyle
\int_{t_0-3.5t_E}^{t_0+3.5t_E}
\left[
A_{\rm BSPL}(t)-A_{\rm PSPL,best}(t)
\right]^2
\,dt
}{
\displaystyle
\int_{t_0-3.5t_E}^{t_0+3.5t_E}
\left[
A_{\rm BSPL}(t)-1
\right]^2
\,dt
}
\right]^{1/2}.
\label{eq:D_metric}
\end{equation}

Thus, $D_{\rm BSPL-PSPL}\rightarrow0$ corresponds to a strong intrinsic
degeneracy, whereas larger values indicate that a larger fraction of
the microlensing signal remains outside the PSPL model family. The
normalization makes $D_{\rm BSPL-PSPL}$ a fractional shape mismatch,
not a survey-specific detection threshold. Because the window is tied
to the microlensing timescale rather than to the binary period, the
long-period limit characterizes similarity over the observable
microlensing event, not over a complete binary orbit. In particular,
for $P\gg t_E$ only a slowly varying fraction of the orbit is present
within the metric window, which is precisely the regime in which the
PSPL fit can absorb much of the perturbation.

\paragraph{Numerical robustness.}
Grid-wide diagnostics, targeted multi-start tests, independent
cross-optimizer refits, temporal-resolution tests, and
integration-window tests confirm that the broad intrinsic structures
used below are stable against the numerical choices of PSPL
initialization, optimizer, temporal sampling, and the adopted
event-centered integration window. All $12{,}000$ fits in the primary
single-luminous-source production grid and all $120{,}600$ fits in the
two-luminous-source grid report successful solutions, with no non-finite
or negative values of $D_{\rm BSPL-PSPL}$. Independent TRF
re-optimization of the same intrinsic objective reproduces the
production Nelder--Mead minima to numerical precision. The full
validation procedure and numerical results are given in
Appendix~\ref{app:numerical_robustness}. Extremely narrow point-source
close approaches require local temporal refinement and are treated
separately in Section~\ref{sec:close_approach}.

The same separation has an observational interpretation. In the
idealized limit of negligible blending, approximately uniform flux
uncertainties, and equivalent temporal weighting,
\begin{equation}
\sqrt{\Delta\chi^2}
\simeq
D_{\rm BSPL-PSPL}\,
{\rm SNR}_{\rm event}.
\end{equation}
For a sampled light curve we define the event signal-to-noise ratio as
\begin{equation}
{\rm SNR}_{\rm event}^{2}
\equiv
\sum_i
\left[
\frac{F_{\rm BSPL}(t_i)-F_{\rm base}}{\sigma_{F,i}}
\right]^2,
\label{eq:event_snr}
\end{equation}
where $F_{\rm base}$ is the unmagnified baseline flux and
$\sigma_{F,i}$ is the flux uncertainty at epoch $t_i$. A real survey
such as Roman will weight the residual through its actual cadence, flux
uncertainties, blending, and observing windows. The corresponding
survey-weighted generalization is defined in
Equation~\ref{eq:D_roman} of Section~\ref{sec:roman_asimov}. We therefore
keep the intrinsic and survey-dependent problems separate.

The second diagnostic is the displacement of the best-fitting PSPL
parameters,
\begin{equation}
\Delta t_0=t_{0,\rm fit}-t_{0,\rm true},
\qquad
\Delta u_0=u_{0,\rm fit}-u_{0,\rm true},
\qquad
\Delta t_E=t_{E,\rm fit}-t_{E,\rm true}.
\end{equation}
These offsets quantify the part of the binary-source perturbation that
is absorbed by moving within the PSPL model family, while
$D_{\rm BSPL-PSPL}$ measures what remains after that optimization.

Figure~\ref{fig:xallarap_examples} illustrates these two regimes for a
dark companion: one configuration leaves a visible irreducible residual,
whereas a long-period configuration is largely absorbed by the PSPL fit.


\begin{figure}[t]
    \centering
    \includegraphics[width=0.49\linewidth]{figures/figures_lightcurves/example_1.png}
    \includegraphics[width=0.49\linewidth]{figures/figures_lightcurves/example_2_confused.png}
    \caption{
    Representative BSPL light curves and their best-fitting PSPL models
    for a dark companion. Left: $a=1\,\mathrm{AU}$ and $P=210$ d,
    leaving a clear irreducible residual. Right: $a=10\,\mathrm{AU}$
    and $P=6000$ d, where the slowly varying perturbation is largely
    absorbed by shifts in $(t_0,u_0,t_E)$. The ``Xallarap'' label
    appearing in the panels refers specifically to this $q_f=0$
    dark-companion limit; elsewhere we use BSPL for the physical
    binary-source system more generally.
    }
    \label{fig:xallarap_examples}
\end{figure}


\section{One Luminous Source}
\label{sec:one_luminous_source}

We first set $q_f=0$, so the companion is dark and the signal is produced
only by the orbital displacement of the luminous source.
Figure~\ref{fig:D_u0_tE} shows $D_{\rm BSPL-PSPL}$ for a
Kepler-consistent family with $M_{S,1}=2\,M_\odot$ and
$M_{S,2}=1\,M_\odot$. We adopt $t_E=30\,$d for the fiducial map because it
lies in the characteristic range of Galactic-bulge microlensing events, while
the adjacent multi-$t_E$ panel retains the longer-duration $t_E=150\,$d case
and demonstrates explicitly how the degeneracy boundary changes with the
absolute event timescale.

The mismatch is small in two different limits. At short periods, Kepler's
law reduces the physical separation and hence $\xi_E$. At long periods,
only a slowly varying part of the orbit is sampled, allowing the PSPL
fit to absorb much of the deformation through $(t_0,u_0,t_E)$. The
largest irreducible residuals occur between these limits. The
quantitative value of the mismatch depends on the projected orbital
geometry, but the broad short--intermediate--long-period structure is
preserved for most orientations tested; a dedicated robustness study is
given in Appendix~\ref{app:geometry_robustness}.

The dependence on $u_0$ shows that $P/t_E$ alone does not organize the
degeneracy. Varying $u_0$ changes both the spatial ratio $\xi_E/u_0$ and
the local temporal ratio $P/(u_0t_E)$. Likewise, at fixed $P/t_E$,
changing $t_E$ changes the physical period and therefore the Keplerian
orbital scale. This is why the $D_{\rm BSPL-PSPL}=10^{-2}$ contours in
Figure~\ref{fig:D_u0_tE} do not collapse onto a single curve.


\begin{figure}[t]
    \centering

    \begin{subfigure}[t]{0.48\linewidth}
        \centering
        \includegraphics[
            width=\linewidth
        ]{../figures/current/D_heatmap_tE30_complete_contour.pdf}

        \caption{
        Normalized BSPL--PSPL mismatch
        $D_{\rm BSPL-PSPL}$ for the fiducial $t_E=30\,{\rm d}$.
        The red contour marks
        $D_{\rm BSPL-PSPL}=10^{-2}$.
        }
        \label{fig:D_map_tE30}
    \end{subfigure}
    \hfill
    \begin{subfigure}[t]{0.48\linewidth}
        \centering
        \includegraphics[
            width=\linewidth
        ]{../figures/current/D_contours_many_tE_complete.pdf}

        \caption{
        Evolution of the
        $D_{\rm BSPL-PSPL}=10^{-2}$ boundary for different
        Einstein timescales.
        }
        \label{fig:D_contours_tE}
    \end{subfigure}

    \caption{
    Intrinsic BSPL--PSPL degeneracy in the single-luminous-source
    limit for the Kepler-consistent binary family considered here.
    Panel~\subref{fig:D_map_tE30} shows
    $D_{\rm BSPL-PSPL}$ in the $(u_0,P/t_E)$ plane for
    $t_E=30\,{\rm d}$. The mismatch becomes small both toward the
    short-period limit, where the Keplerian orbital amplitude decreases,
    and toward the long-period limit, where the slowly varying orbital
    perturbation can be increasingly absorbed by a re-optimized PSPL
    trajectory.
    Panel~\subref{fig:D_contours_tE} shows the
    $D_{\rm BSPL-PSPL}=10^{-2}$ boundary for different $t_E$.
    The displacement of the contours demonstrates that
    $(u_0,P/t_E)$ alone does not uniquely specify the degeneracy in a
    Kepler-consistent physical family.
    The level $D_{\rm BSPL-PSPL}=10^{-2}$ is used only as a reference
    fractional mismatch, not as a survey-specific detection threshold.
    }
    \label{fig:D_u0_tE}
\end{figure}

\subsection{Absorption by the PSPL fit}
\label{sec:pspl_biases}

A small residual does not imply unbiased PSPL parameters.
Figure~\ref{fig:tE_bias_D} compares the relative $t_E$ bias with
contours of $D_{\rm BSPL-PSPL}$. The bias changes magnitude and sign
without following the residual contours, demonstrating that the absorbed
and irreducible components of the perturbation carry different
information.

The mass-ratio scan provides a useful limiting test. At fixed total
source mass,
\begin{equation}
\xi_{E,1}
=
\frac{q_M}{1+q_M}\,\xi_{\rm rel}
\longrightarrow0
\qquad
(q_M\rightarrow0),
\end{equation}
so both the mismatch and the PSPL parameter offsets vanish toward the
single-source limit.


\begin{figure}[t]
    \centering
    \includegraphics[width=1\linewidth]
    {../figures/current/paper_tE_bias_D_contours_tE30.pdf}

    \caption{
    Relative bias in the Einstein timescale recovered by the best-fitting
    PSPL model,
    $\Delta t_E/t_E=
    (t_{E,\mathrm{fit}}-t_{E,\mathrm{true}})/t_{E,\mathrm{true}}$.
    Black curves show contours of the normalized irreducible mismatch
    $D_{\rm BSPL-PSPL}$.
    Panel~(a) shows the scan over $u_0$ at fixed
    $q_M=0.5$, while panel~(b) shows the scan over $q_M$ at fixed
    $u_0=0.1$ and fixed total source mass.
    The two panels use independent symmetric color ranges for
    $\Delta t_E/t_E$ because the robust bias amplitudes differ between
    the $u_0$ and $q_M$ scans. The different structures of the bias and
    the residual mismatch show that the displacement of the best-fitting
    PSPL parameters and the irreducible BSPL residual probe complementary
    components of the degeneracy.
    }
    \label{fig:tE_bias_D}
\end{figure}

The fitted $u_0$ and $t_E$ offsets are also strongly coupled.
For the fiducial $t_E=30\,$d slice at $u_0=0.09885$ and $q_M=0.5$,
Figure~\ref{fig:u0_tE_bias_slice} gives a Spearman coefficient
$\rho_s=0.98666$ ($p=1.98\times10^{-47}$). The PSPL fit therefore
absorbs source orbital motion predominantly along a correlated direction
in $(u_0,t_E)$ rather than through independent parameter shifts.


\begin{figure}[t]
    \centering
    \includegraphics[width=0.5\linewidth]
    {../figures/current/paper_u0_tE_correlated_bias_tE30.pdf}

    \caption{
    Relative shifts in the fitted PSPL impact parameter and Einstein
    timescale as a function of $P/t_E$ for
    $u_0=0.1$ and $q_M=0.5$.
    For the regenerated $t_E=30\,$d products, the plotted slice uses
    $u_0=0.09885$ and $q_M=0.5$, giving a Spearman correlation
    $\rho_s=0.98666$ with $p=1.98\times10^{-47}$.
    }
    \label{fig:u0_tE_bias_slice}
\end{figure}

These results establish the central distinction used throughout this
work: a BSPL event can lie close to the PSPL model family while the
recovered PSPL parameters remain systematically displaced.


\section{Two Luminous Sources}
\label{sec:two_luminous_sources}

When both stars contribute light, the observed magnification is
\begin{equation}
A_{\rm BSPL}(t)
=
\frac{A_1(t)+q_fA_2(t)}{1+q_f},
\qquad
q_f=\frac{F_{S,2}}{F_{S,1}}.
\end{equation}
Figure~\ref{fig:qmass_qflux_D} shows that the $(q_M,q_f)$ plane contains
a broad region of enhanced degeneracy near $q_f=q_M$.

The origin is the motion of the unresolved source photocenter. For
relative orbital displacement $\xi_{\rm rel}\mathbf{s}(t)$,
\begin{equation}
\delta\mathbf{u}_1
=
\frac{q_M}{1+q_M}\xi_{\rm rel}\mathbf{s}(t),
\qquad
\delta\mathbf{u}_2
=
-\frac{1}{1+q_M}\xi_{\rm rel}\mathbf{s}(t),
\end{equation}
so the flux-weighted displacement is
\begin{equation}
\delta\mathbf{u}_{\rm phot}
=
\frac{q_M-q_f}{(1+q_M)(1+q_f)}
\xi_{\rm rel}\mathbf{s}(t).
\label{eq:photocenter_displacement}
\end{equation}
The first-order photocenter motion therefore vanishes for $q_f=q_M$,
suppressing the leading orbital perturbation.


\begin{figure}[t]
    \centering

    \includegraphics[
        width=\linewidth
    % Source product: results/final_483d90a0fd07/qmass_qflux_tE30/summary_qM_qf.npz
    ]{../figures/current/qmass_qflux_D_map.pdf}

    \caption{
    Intrinsic BSPL--PSPL mismatch in the
    $(q_M,q_f)$ plane for the fiducial $t_E=30\,$d and
    $P/t_E=0.3$, $1$, and $3$ (left to right), corresponding to
    physical periods $P=9$, $30$, and $90\,$d.
    The color scale shows
    $D_{\rm BSPL-PSPL}$ after refitting each binary-source
    light curve with a PSPL model.
    The dashed diagonal marks $q_f=q_M$, where the
    first-order displacement of the unresolved source
    photocenter vanishes.
    Localized narrow structures away from this locus can arise from
    extreme source--lens close approaches and can be sensitive to the
    temporal resolution of a uniform point-source calculation. They are
    discussed in Section~\ref{sec:close_approach} and are not associated
    with the photocenter-cancellation mechanism.
    }

    \label{fig:qmass_qflux_D}
\end{figure}

\subsection{Photocenter cancellation}
\label{sec:photocenter_cancellation}

The cancellation does not make the BSPL light curve exactly PSPL because
higher-order terms remain. This can be seen directly by expanding the
combined magnification around the barycentric trajectory. Let
$\boldsymbol{\delta}=\xi_{\rm rel}\mathbf{s}(t)$,
$a=q_M/(1+q_M)$, $b=1/(1+q_M)$, and
$A_B=A(\mathbf{u}_B)$. To second order,
\begin{equation}
A(\mathbf{u}_B+\delta\mathbf{u})
=
A_B
+
\nabla A_B\cdot\delta\mathbf{u}
+
\frac{1}{2}
\delta\mathbf{u}^{\rm T}\mathbf{H}_A\delta\mathbf{u}
+
\mathcal{O}(\delta u^3),
\label{eq:taylor_A}
\end{equation}
where $\mathbf{H}_A$ is the Hessian of the point-lens magnification
evaluated on the barycentric trajectory. Using
$\delta\mathbf{u}_1=a\boldsymbol{\delta}$ and
$\delta\mathbf{u}_2=-b\boldsymbol{\delta}$ in the flux-weighted
binary-source magnification gives
\begin{align}
A_{\rm BSPL}-A_B
={}&
\frac{q_M-q_f}{(1+q_M)(1+q_f)}
\nabla A_B\cdot\boldsymbol{\delta}
\nonumber\\
&+
\frac{q_M^2+q_f}
{2(1+q_M)^2(1+q_f)}
\boldsymbol{\delta}^{\rm T}\mathbf{H}_A\boldsymbol{\delta}
+
\mathcal{O}(\xi_{\rm rel}^3).
\label{eq:photocenter_taylor}
\end{align}
Thus the linear term vanishes identically on $q_f=q_M$, while the
quadratic term remains,
\begin{equation}
A_{\rm BSPL}-A_B
=
\frac{q_M}{2(1+q_M)^2}
\xi_{\rm rel}^{2}
\mathbf{s}^{\rm T}\mathbf{H}_A\mathbf{s}
+
\mathcal{O}(\xi_{\rm rel}^{3})
\qquad(q_f=q_M).
\label{eq:photocenter_quadratic}
\end{equation}
The PSPL re-optimization can absorb components of these deformations,
but it cannot reintroduce a first-order term that is absent from the
BSPL light curve itself. We therefore expect a generic linear mismatch
away from the cancellation locus and a quadratic leading scaling on the
locus.

To test this prediction numerically after projection onto the best-fitting
PSPL family, we vary $\xi_{\rm rel}$ independently of the Keplerian
relation between period and orbital size. Figure~\ref{fig:photocenter_scaling} shows
\begin{equation}
D_{\rm BSPL-PSPL}\propto\xi_{\rm rel}
\qquad (q_f=0),
\end{equation}
whereas along the cancellation condition
\begin{equation}
D_{\rm BSPL-PSPL}\propto\xi_{\rm rel}^{2}
\qquad (q_f=q_M).
\end{equation}
The fitted logarithmic slopes are numerically consistent with one and
two, respectively, in agreement with the first- and second-order terms
in the expansion above. The dedicated orientation test in
Appendix~\ref{app:geometry_robustness} gives median slopes
$\gamma_{\rm dark}=1.0000$ and
$\gamma_{\rm cancel}=2.0000$; the full orientation-to-orientation
ranges are reported there. The
numerical experiment therefore verifies, after PSPL re-optimization,
the analytic expectation that photocenter cancellation removes the
first-order binary-source deformation rather than the binary-source
signal itself. The linear-to-quadratic change in small-separation scaling persists
across the projected orbital orientations explored there and is
therefore not a consequence of the fiducial
$(\alpha,\phi,i)=(0,0,\pi/2)$ geometry.


\begin{figure}[t]
    \centering

    \includegraphics[
        width=\linewidth
    ]{../figures/current/photocenter_scaling.pdf}

    \caption{
    Small-separation scaling of the intrinsic
    BSPL--PSPL mismatch.
    For a dark companion ($q_f=0$), the residual mismatch
    scales linearly with the relative orbital amplitude,
    $D_{\rm BSPL-PSPL}\propto\xi_{\rm rel}$.
    Along the photocenter-cancellation condition
    $q_f=q_M$, the linear contribution vanishes and the
    mismatch becomes quadratic,
    $D_{\rm BSPL-PSPL}\propto\xi_{\rm rel}^2$.
    The right panel shows the corresponding fitted
    logarithmic slopes.
    }

    \label{fig:photocenter_scaling}
\end{figure}

\subsection{Astrophysical accessibility of photocenter suppression}
\label{sec:physical_photocenter}

The condition $q_f=q_M$ is an exact geometric cancellation locus, not a
generic mass--luminosity relation for main-sequence binaries. This
distinction is important for interpreting Figure~\ref{fig:qmass_qflux_D}.
Empirical main-sequence mass--luminosity relations are strongly
mass-dependent rather than described by a single power law
\citep{Eker2018}. We therefore use a stellar-physics calculation to ask a
more limited and directly relevant question: can physically plausible
main-sequence binaries approach the cancellation locus closely enough to
suppress the first-order photocenter motion?

We use a 10-Gyr, solar-metallicity PARSEC v1.2S isochrone
\citep{Bressan2012}, selecting only main-sequence points and synthetic
magnitudes in the Roman 2024 F146 passband. A similar use of a 10-Gyr,
solar-metallicity PARSEC isochrone to infer binary-source flux ratios has
been adopted in detailed microlensing analyses \citep{Miyazaki2020}. For
primary masses $M_1=0.5$, $0.7$, and $0.8\,M_\odot$, all comfortably below
the turnoff of the adopted isochrone, we set $M_2=q_M M_1$ and compute
\begin{equation}
q_f^{\rm F146}
=
10^{-0.4\left(M_{{\rm F146},2}-M_{{\rm F146},1}\right)}.
\label{eq:qf_f146}
\end{equation}
The F146 calculation is used only to test the physical accessibility of
the photocenter mechanism. It is distinct from the legacy W149
photometric-error prescription used in the Roman Asimov experiment of
Section~\ref{sec:roman_asimov}.

A useful dimensionless measure is the surviving first-order photocenter
amplitude relative to the same binary with a dark companion. From
Equation~\ref{eq:photocenter_displacement},
\begin{equation}
S_{\rm ph}
\equiv
\frac{|\delta u_{\rm phot}|}
{\left.|\delta u_{\rm phot}|\right|_{q_f=0}}
=
\frac{|q_M-q_f|}{q_M(1+q_f)}.
\label{eq:S_ph}
\end{equation}
Thus $S_{\rm ph}=1$ corresponds to no luminous-companion suppression and
$S_{\rm ph}=0$ to exact first-order cancellation.
Figure~\ref{fig:physical_photocenter} shows that the PARSEC F146 tracks do
not follow $q_f=q_M$ except in the equal-mass limit. Nevertheless, strong
partial suppression is physically accessible near $q_M=1$. At
$q_M=0.95$, the three adopted primary masses give
$S_{\rm ph}=0.064$, $0.080$, and $0.089$, respectively, so only about
$6$--$9\%$ of the dark-companion first-order photocenter amplitude
survives. At $q_M=0.8$, the surviving fraction is approximately
$0.24$--$0.37$.

\begin{figure}[t]
    \centering

    \begin{subfigure}[t]{0.49\linewidth}
        \centering
        \includegraphics[width=\linewidth]
        {../figures/current/physical_qM_qf_PARSECF146.pdf}
        \caption{
        Roman F146 flux-ratio tracks compared with the exact
        photocenter-cancellation locus $q_f=q_M$ and the illustrative
        toy relation $q_f=q_M^4$.
        }
    \end{subfigure}
    \hfill
    \begin{subfigure}[t]{0.49\linewidth}
        \centering
        \includegraphics[width=\linewidth]
        {../figures/current/physical_photocenter_suppression_PARSECF146.pdf}
        \caption{
        Surviving first-order photocenter amplitude $S_{\rm ph}$ relative
        to the corresponding dark-companion binary.
        }
    \end{subfigure}

    \caption{
    Astrophysical accessibility of photocenter suppression for a
    representative old main-sequence population. The tracks are derived
    from a 10-Gyr, solar-metallicity PARSEC isochrone with Roman F146
    synthetic photometry for $M_1=0.5$, $0.7$, and $0.8\,M_\odot$.
    Exact cancellation is not generic for unequal-mass main-sequence
    binaries, but nearly equal-mass systems can strongly suppress the
    first-order photocenter motion. This figure tests the physical
    accessibility of the analytic mechanism; it is not a population
    synthesis and is not used to assign population weights to the
    $(q_M,q_f)$ microlensing grid.
    }
    \label{fig:physical_photocenter}
\end{figure}

The physical isochrone test and the fixed-$M_{\rm tot}$ microlensing grid
answer complementary questions. The former determines which flux ratios
are plausible for stellar binaries, while the latter isolates the
geometric dependence of the BSPL--PSPL projection at fixed orbital scale.
Because the total source mass varies along an isochrone track, we do not
interpret an interpolation through the fixed-$M_{\rm tot}=3\,M_\odot$
map as a self-consistent prediction of $D_{\rm BSPL-PSPL}$ for a stellar
population. A population-level calculation would additionally require
stellar ages and metallicities, binary demographics, Galactic weights,
and the survey selection function.

The same configurations can remain close to PSPL while shifting the
inferred microlensing parameters. Figure~\ref{fig:qmass_qflux_biases}
shows that the parameter-bias structure is complementary to
$D_{\rm BSPL-PSPL}$, as in the single-luminous-source case.


\begin{figure}[t]
    \centering

    \includegraphics[
        width=\linewidth
    % Source product: results/final_483d90a0fd07/qmass_qflux_tE30/summary_qM_qf.npz
    ]{../figures/current/qmass_qflux_biases.pdf}

    \caption{
    Relative shifts of the best-fitting PSPL parameters
    across the two-luminous-source $(q_M,q_f)$ experiment for the
    fiducial $t_E=30\,$d grid.
    The parameter biases measure the component of the
    binary-source perturbation absorbed by displacement
    within the PSPL model family and are complementary to
    the irreducible mismatch $D_{\rm BSPL-PSPL}$.
    }

    \label{fig:qmass_qflux_biases}
\end{figure}

\subsection{Localized close-approach features}
\label{sec:close_approach}

The point-source maps contain isolated narrow structures produced by
extreme source--lens close approaches. Such localized features in the
$P/t_E=1$ panel of Figure~\ref{fig:qmass_qflux_D} belong to this class
rather than to the photocenter-cancellation ridge. Their precise location
can shift when the absolute timescale is changed because the Keplerian
orbital scale changes at fixed $P/t_E$. Targeted calculations show that such
features can be unresolved by a uniform $10^4$-epoch grid and can also
be substantially smoothed by finite-source effects. They are therefore
not used to infer the broad BSPL--PSPL degeneracy structure or the
photocenter-cancellation mechanism. The temporal-resolution,
local-refinement, and finite-source checks are given in
Appendix~\ref{app:numerical_robustness}.

\section{Roman Asimov model separation}
\label{sec:roman_asimov}

The intrinsic metric in Equation~\ref{eq:D_metric} deliberately removes
survey-specific weighting. To connect the same physical configurations
to a survey-like experiment, we propagate the fiducial $t_E=30\,$d
single-luminous-source ($q_f=0$) grid through noiseless Roman W149 Asimov
light curves. The Roman calculation uses exactly the same 133
$u_0\leq1$ nodes and 60 periods as the corresponding intrinsic grid and
three source magnitudes, W149$=19$, 21, and 23. The matched production
therefore contains $133\times60\times3=23{,}940$ PSPL fits, all of which
were completed in the production run.

The cadence and photometric-precision model are a frozen Roman-like
prescription inherited from the simulation pipeline rather than a current
operations forecast. The event is centered on the third nominal
high-cadence season. For $t_E=30\,$d, the event-centered interval
$|t-t_0|\leq3.5t_E$ contains 8,688 W149 epochs. Table~\ref{tab:roman_setup}
summarizes the ingredients that enter the matched calculation. The W149
nomenclature is retained here because it identifies the legacy
photometric prescription used in the production pipeline.

\begin{table}[H]
\centering
\small
\caption{
Roman-like W149 Asimov setup used for the matched $t_E=30\,$d
survey-weighted experiment. The cadence and uncertainty prescription are
frozen inputs to the calculation and are not intended as a current Roman
operations schedule.
}
\label{tab:roman_setup}
\begin{tabular}{p{0.34\linewidth}p{0.56\linewidth}}
\toprule
Quantity & Adopted prescription \\
\midrule
Einstein timescale & $t_E=30\,$d \\
Intrinsic nodes retained & 133 values of $u_0\leq1$ and 60 periods \\
Source magnitudes & W149$=19$, 21, 23 \\
Total PSPL fits & 23,940 \\
Event placement & $t_0$ at the midpoint of the third nominal high-cadence season \\
Analysis window & $|t-t_0|\leq3.5t_E$ \\
Retained W149 epochs & 8,688 \\
High-cadence sampling & 12.1 min within the adopted active season \\
W149 flux zero point & 27.615 \\
Photometric uncertainty & Log-linear interpolation of the adopted W149 magnitude--precision relation, with a bright-end floor of $10^{-3}$ mag \\
True blending & No additional blended flux in the BSPL Asimov truth ($F_b=0$) \\
PSPL optimizer & pyLIMA trust-region reflective fit initialized at the true $(t_0,u_0,t_E)$ \\
PSPL nonlinear bounds & $t_0\in[t_{0,\rm true}-2t_E,t_{0,\rm true}+2t_E]$, $u_0\in[-20,20]$, and $t_E\in[0.02t_{E,\rm true},20t_{E,\rm true}]$ \\
Photometric nuisance parameters & Source flux and total baseline flux re-optimized with the \texttt{ftotal} parameterization; no $F_b\geq0$ prior is imposed \\
\bottomrule
\end{tabular}
\end{table}

Allowing formally negative fitted blend flux enlarges the PSPL nuisance
parameter space. Relative to a fit with a non-negative blend constraint,
this can only reduce or leave unchanged the minimized PSPL $\chi^2$ and
therefore gives a conservative model separation for the adopted Asimov
light curves.

For each configuration, the Asimov fluxes are generated from the BSPL
truth without a random noise realization, while the W149 uncertainties
are retained as statistical weights. We refit a PSPL model in flux space,
allowing $(t_0,u_0,t_E)$ and the photometric source and total-flux nuisance
parameters to re-optimize. The resulting Asimov model-separation statistic
is
\begin{equation}
\Delta\chi^2_{\rm Roman}
=
\chi^2_{\rm PSPL,best}-\chi^2_{\rm BSPL}
=
\sum_i
\left[
\frac{F_{\rm BSPL}(t_i)-F_{\rm PSPL,best}(t_i)}{\sigma_{F,i}}
\right]^2,
\label{eq:roman_dchi2}
\end{equation}
where the final equality follows because the noiseless BSPL Asimov data
have $\chi^2_{\rm BSPL}=0$. We also define the survey-weighted normalized
mismatch
\begin{equation}
D_{\rm Roman,w}
\equiv
\frac{\sqrt{\Delta\chi^2_{\rm Roman}}}{{\rm SNR}_{\rm event}}
=
\left[
\frac{
\displaystyle\sum_i
\left[
\frac{F_{\rm BSPL}(t_i)-F_{\rm PSPL,best}(t_i)}{\sigma_{F,i}}
\right]^2
}{
\displaystyle\sum_i
\left[
\frac{F_{\rm BSPL}(t_i)-F_{\rm base}}{\sigma_{F,i}}
\right]^2
}
\right]^{1/2}.
\label{eq:D_roman}
\end{equation}
This quantity is not identical to $D_{\rm BSPL-PSPL}$ because the Roman
fit changes both the temporal weighting and the fitting objective, and
also re-optimizes photometric nuisance parameters.

Figure~\ref{fig:roman_intrinsic_comparison} provides the direct matched
comparison. The heat map shows $\log_{10}\Delta\chi^2_{\rm Roman}$,
while the intrinsic $D_{\rm BSPL-PSPL}=10^{-2}$ and $10^{-3}$ contours
are overlaid on the same $(u_0,P/t_E)$ nodes. The broad survey-weighted
gradient follows the intrinsic topology, but the reference
$\Delta\chi^2_{\rm Roman}=100$ contour cuts across the intrinsic levels
and moves systematically toward smaller $P/t_E$ as the source becomes
fainter. Brighter sources can therefore distinguish some configurations
that are already close to the PSPL family in the intrinsic metric, while
fainter sources require a larger residual mismatch to reach the same
expected $\Delta\chi^2$.

\begin{figure}[t]
    \centering
    \includegraphics[width=\linewidth]
    {../figures/current/roman_intrinsic_comparison.pdf}
    \caption{
    Matched $t_E=30\,$d Roman-like Asimov model separation for the
    single-luminous-source grid. The color scale shows
    $\log_{10}\Delta\chi^2_{\rm Roman}$ for W149$=19$, 21, and 23.
    The solid and dashed intrinsic contours mark
    $D_{\rm BSPL-PSPL}=10^{-2}$ and $10^{-3}$, respectively. The white
    contour marks $\Delta\chi^2_{\rm Roman}=100$, used only as a common
    visual reference. Its displacement with source magnitude illustrates
    that an intrinsic model-space distance does not map to a universal
    survey threshold.
    }
    \label{fig:roman_intrinsic_comparison}
\end{figure}

The grid is logarithmically sampled in $(u_0,P)$ and is not weighted by
the source-binary population, Galactic event rate, or Roman event
selection. We therefore do not convert the area of the plotted regions
into an expected fraction of Roman binary-source events. Likewise,
$\Delta\chi^2_{\rm Roman}=100$ is not assigned a $p$-value. A threshold
with a specified false-positive rate and statistical power requires the
null distribution of a likelihood-ratio statistic to be calibrated with
noisy realizations and both competing fits. The Asimov calculation here
has the narrower purpose of translating the intrinsic model-space
structure into the expected survey-weighted residual power for a fixed
observing prescription.

\section{Conclusions}
\label{sec:conclusions}

We have characterized how binary-source microlensing light curves
project onto the PSPL model family after $(t_0,u_0,t_E)$ are re-optimized.
The normalized distance $D_{\rm BSPL-PSPL}$ measures the irreducible
component of the binary-source signal, while the fitted parameter offsets
measure the component absorbed by motion within the PSPL family.

For one luminous source, the degeneracy is controlled by both spatial
and temporal scales. In the fiducial $t_E=30\,$d grid the mismatch is
small at short periods, where the Keplerian orbital displacement becomes
small, and again at long periods, where only a slowly varying fraction of
the orbit is sampled and the PSPL fit can absorb much of the deformation.
The boundary therefore cannot be described by $P/t_E$ alone. Moreover, a
small residual does not imply unbiased PSPL parameters: at
$u_0=0.09885$ and $q_M=0.5$, the relative $u_0$ and $t_E$ shifts have
$\rho_s=0.98666$.

For two luminous sources, the analytic condition $q_f=q_M$ removes the
first-order displacement of the unresolved photocenter. The resulting
small-separation mismatch changes from
$D_{\rm BSPL-PSPL}=\mathcal{O}(\xi_{\rm rel})$ for a dark companion to
$D_{\rm BSPL-PSPL}=\mathcal{O}(\xi_{\rm rel}^2)$ on the cancellation
locus. The orientation study shows that this change in asymptotic order
is independent of projected orbital geometry. The cancellation locus is
not, however, a generic main-sequence mass--flux relation. Using a
representative 10-Gyr solar-metallicity PARSEC isochrone in Roman F146,
we find instead that its astrophysical relevance is concentrated toward
nearly equal-mass binaries: for $q_M=0.95$, only about $6$--$9\%$ of the
first-order dark-companion photocenter amplitude survives for the three
primary masses tested. Thus the physically relevant result is strong
\emph{partial} photocenter suppression near the equal-mass limit, with
exact cancellation providing the analytic limiting case.

The matched Roman-like Asimov experiment translates the same
$t_E=30\,$d intrinsic grid into survey-weighted residual power. The
production contains 23,940 PSPL fits at W149$=19$, 21, and 23. The
survey-weighted maps broadly preserve the intrinsic degeneracy topology,
but the $\Delta\chi^2_{\rm Roman}=100$ reference contour shifts with
source brightness and cuts across fixed intrinsic-$D$ contours. The
intrinsic distance from the PSPL model family is therefore a physical
property of the light-curve morphology, whereas the expected survey
separation also depends on cadence and photometric precision.

Targeted numerical tests support the broad structures used in this
interpretation. Multi-start and independent cross-optimizer fits recover
the production intrinsic minima, reasonable changes in the integration
window preserve the short--intermediate--long-period ordering, and the
linear-to-quadratic photocenter scaling survives changes in projected
orbital geometry. Extremely narrow point-source close approaches are a
separate regime: they can require local temporal refinement and can be
substantially smoothed by finite-source effects, so they are not used to
define the generic degeneracy boundaries.

The present analysis intentionally isolates a restricted model space.
The lens is point-like, the source orbit is circular, and the intrinsic
experiments omit annual parallax, orbital eccentricity, and other
higher-order effects. The stellar-isocrone calculation is a physical
sanity check rather than a population synthesis, and the Roman Asimov
grid is not weighted by binary demographics or the Galactic event rate.
Calibrating a binary-source detection threshold with a specified
false-positive rate will require noisy likelihood-ratio simulations;
turning the present maps into an expected Roman event fraction will
additionally require population and selection-function weights. These
are distinct extensions of the model-space degeneracy characterized
here.

Overall, the results separate three quantities that should not be
conflated: the intrinsic distance of a BSPL light curve from the PSPL
model family, the bias absorbed into the best-fitting PSPL parameters,
and the survey-weighted significance of the remaining residual. This
separation provides a compact framework for identifying when source
multiplicity can remain hidden and for determining which additional
physical or observational information is required to reveal it.

\appendix
\section{Binary-source orbital-motion model}
\label{app:bspl_formalism}

Consider source masses $M_{S,1}$ and $M_{S,2}$ with
\begin{equation}
q_M=\frac{M_{S,2}}{M_{S,1}}.
\end{equation}
If $a_{\rm rel}$ is the semimajor axis of the relative orbit, the
individual barycentric semimajor axes are
\begin{equation}
a_1=\frac{q_M}{1+q_M}a_{\rm rel},
\qquad
a_2=\frac{1}{1+q_M}a_{\rm rel},
\end{equation}
with
\begin{equation}
a_{\rm rel}^3
=
\frac{G(M_{S,1}+M_{S,2})P^2}{4\pi^2}.
\end{equation}
The corresponding dimensionless amplitudes are
\begin{equation}
\xi_{\rm rel}=\frac{a_{\rm rel}}{\hat r_E},
\qquad
\xi_{E,1}=\frac{q_M}{1+q_M}\xi_{\rm rel},
\qquad
\xi_{E,2}=\frac{1}{1+q_M}\xi_{\rm rel},
\end{equation}
where $\hat r_E=D_S\theta_E$.

For a circular orbit we write
\begin{equation}
\mathbf{s}(t)
=
\begin{pmatrix}
\cos\Omega(t)\\
\cos i\,\sin\Omega(t)
\end{pmatrix},
\qquad
\Omega(t)
=
\frac{2\pi}{P}(t-t_{\rm ref})+\phi.
\end{equation}
After rotation by the orientation angle $\alpha$,
\begin{equation}
\mathbf{s}_{\rm proj}(t)
=
\mathbf{R}(\alpha)\mathbf{s}(t),
\qquad
\mathbf{R}(\alpha)
=
\begin{pmatrix}
\cos\alpha & -\sin\alpha\\
\sin\alpha & \cos\alpha
\end{pmatrix}.
\end{equation}
The source displacements relative to the barycenter are
\begin{equation}
\boldsymbol{\delta u}_1(t)
=
\frac{q_M}{1+q_M}\xi_{\rm rel}\mathbf{s}_{\rm proj}(t),
\qquad
\boldsymbol{\delta u}_2(t)
=
-\frac{1}{1+q_M}\xi_{\rm rel}\mathbf{s}_{\rm proj}(t).
\end{equation}

The rectilinear barycentric trajectory is
\begin{equation}
\boldsymbol{u}_{\rm B}(t)
=
\begin{pmatrix}
(t-t_0)/t_E\\
u_0
\end{pmatrix},
\end{equation}
and the component trajectories are
\begin{equation}
\boldsymbol{u}_i(t)
=
\boldsymbol{u}_{\rm B}(t)
+
\boldsymbol{\delta u}_i(t),
\qquad i=1,2.
\end{equation}
For $u_i(t)=|\boldsymbol{u}_i(t)|$, each source has PSPL magnification
\begin{equation}
A_i(t)
=
\frac{u_i^2(t)+2}
{u_i(t)\sqrt{u_i^2(t)+4}}.
\end{equation}

If the unmagnified source fluxes are $F_{S,1}$ and $F_{S,2}$, define
\begin{equation}
q_f=\frac{F_{S,2}}{F_{S,1}}.
\end{equation}
Ignoring additional blended light in the intrinsic calculations, the
combined binary-source magnification normalized to the total source flux
is
\begin{equation}
A_{\rm BSPL}(t)
=
\frac{A_1(t)+q_fA_2(t)}{1+q_f}.
\end{equation}
The limit $q_f\rightarrow0$ recovers the single-luminous-source case.




\section{Numerical robustness of the intrinsic projection}
\label{app:numerical_robustness}

For reference, Table~\ref{tab:numerical_robustness_summary} summarizes the
numerical checks collected in this appendix. The tests are deliberately
complementary: some probe sensitivity to initialization or optimizer,
whereas others test grid completeness, temporal sampling, fitting window,
finite-source smoothing, or imposed parameter bounds.

\begin{table*}[t]
\centering
\caption{Summary of the numerical robustness checks used in this work.}
\label{tab:numerical_robustness_summary}
\begin{tabular}{p{0.22\textwidth}p{0.43\textwidth}p{0.27\textwidth}}
\toprule
Check & Method & Main result \\
\midrule
Multi-start convergence & Five representative configurations, 22 initial conditions each, including 12 broad random starts & No scientifically distinct lower minimum; largest fractional change $8\times10^{-14}$ \\
One-luminous grid audit & All $12{,}000$ stored production fits & $12{,}000/12{,}000$ successful; no non-finite or negative $D$ \\
Two-luminous grid audit & All $120{,}600$ stored $(q_M,q_f,P/t_E)$ fits & $120{,}600/120{,}600$ successful; no non-finite $D$ \\
Optimizer cross-check & 30 configurations, Nelder--Mead versus bounded TRF on the identical intrinsic objective & Maximum fractional difference in $D$: $1.4\times10^{-13}$; no active TRF bounds \\
Stratified re-fit audit & 60 production points spanning $D$ quantiles, $D\simeq10^{-2}$, parameter-bias extrema, and $(u_0,P)$ boundaries; each re-fitted with fresh Nelder--Mead and same-objective TRF & Stored minima reproduced to numerical precision \\
Temporal resolution & $5\times10^3$--$8\times10^4$ uniform epochs for representative broad features & $10^4$-epoch production values stable to better than $10^{-12}$ fractionally \\
Localized close approach & Dense local temporal refinement of an extreme point-source excursion & Converged to $D=0.569329$; uniform grids can under-resolve the feature \\
Finite-source sanity check & Finite-source truth with varied $\rho_2$ for control and extreme close-approach cases & Control unchanged; extreme mismatch decreases monotonically from $0.569329$ to $0.086$, implying a shift toward greater degeneracy \\
Integration window & Re-fits over $t_0\pm2.5t_E$, $\pm3.5t_E$, and $\pm5t_E$ & Broad features stable; largest changes are percent-level in slowly varying long-period cases \\
Roman $t_E$ bounds & Grid-wide $1\%$ proximity audit plus upper-bound relaxation from $20$ to $200\,t_{E,\rm true}$ & No lower-bound cases; one upper-bound configuration persists across all magnitudes, with negligible effect on model separation \\
\bottomrule
\end{tabular}
\end{table*}

The production intrinsic PSPL fits use Nelder--Mead initialized at the
true barycentric $(t_0,u_0,t_E)$. Because this is a local optimizer, we
tested whether substantially different initial conditions could reach a
lower value of the intrinsic objective. Five representative
configurations were selected to span the strongly degenerate
long-period regime, the smallest impact parameter in the fiducial grid,
the photocenter-cancellation locus, and ordinary and extreme
close-approach geometries. For each configuration we repeated the fit
from 22 alternative initial conditions, including 12 broad random
starts. None of these tests produced a scientifically distinct lower
minimum than the production solution. The largest fractional difference
between the best multi-start result and the truth-initialized fit was
$8\times10^{-14}$, consistent with floating-point-level numerical
differences. Some distant starts instead converged to substantially
poorer local minima, confirming that the objective is not globally
convex even though the production solution is stable for the tested
configurations.

We complemented these targeted tests with a grid-wide audit of the
stored production results. All $12{,}000/12{,}000$ fits in the primary
single-luminous-source $(u_0,P)$ grid have \texttt{SUCCESS=True}, with
no non-finite or negative values of $D_{\rm BSPL-PSPL}$. The largest
change in $\log_{10}D_{\rm BSPL-PSPL}$ between adjacent cells is only
$0.114$ dex. It occurs along the period direction at
$P/t_E=21.50$--$25.13$, where $D_{\rm BSPL-PSPL}$ changes from
$7.32\times10^{-3}$ to $5.63\times10^{-3}$. This diagnostic belongs
to the one-luminous $(u_0,P)$ grid and is therefore unrelated to the
cusp-like structure near $q_M\simeq0.1$ at $P/t_E=1$ discussed in
Section~\ref{sec:close_approach}, which arises in the distinct
two-luminous $(q_M,q_f)$ parameter space. The two-luminous-source grid
likewise contains $120{,}600/120{,}600$ successful fits and no
non-finite values of $D_{\rm BSPL-PSPL}$. These diagnostics show no
grid-wide population of failed or obviously pathological solutions.

To separate optimizer dependence from objective-function dependence, we
also re-optimized exactly the same intrinsic magnification-space
objective with a trust-region reflective solver. For this cross-check we
used $t_0\in[t_{0,\rm ref}-2t_{E,\rm ref},\,
 t_{0,\rm ref}+2t_{E,\rm ref}]$, $u_0\in[-20,20]$, and
$t_E\in[0.02,20]\,t_{E,\rm ref}$. Across 30 systematically selected
configurations, the TRF and production Nelder--Mead values of
$D_{\rm BSPL-PSPL}$ agree to numerical precision: the maximum absolute
fractional difference is $1.4\times10^{-13}$ and the median is
$3.3\times10^{-15}$. None of the TRF solutions reached an active
parameter bound.

As a separate production-grid audit, we selected 60 configurations
stratified across the full $D_{\rm BSPL-PSPL}$ distribution. The
selection includes mismatch quantiles, points close to
$D_{\rm BSPL-PSPL}=10^{-5},10^{-4},10^{-3},10^{-2},10^{-1}$ and $0.3$,
the largest fitted-parameter excursions, and the boundaries of the
sampled $(u_0,P)$ domain. Each selected configuration was independently
re-fitted twice: first with a fresh Nelder--Mead minimization matching
the production procedure, and then with the bounded TRF solver applied
to the same intrinsic magnification-space objective. Both solutions
were compared with each other and with the stored production value.
Figure~\ref{fig:stratified_refit_audit} summarizes this comparison; the
stored minima are reproduced to numerical precision throughout the
stratified sample.

Thus the difference between the intrinsic and Roman metrics is not
produced by the numerical optimizer: it arises when the objective is
changed from the intrinsic magnification-space projection to the
photometrically weighted Roman flux-space projection, in which the flux
nuisance parameters are also re-optimized.

\begin{figure*}[t]
    \centering
    \includegraphics[width=0.92\textwidth]{../figures/appendix/optimizer_stratified_refit_audit.pdf}
    \caption{
    Independent re-fit audit of 60 stratified points from the
    one-luminous production grid. Panel (a) compares the mismatch
    obtained from a fresh Nelder--Mead fit with the value stored in the
    production grid; the diagonal denotes equality. Panel (b) shows the
    absolute fractional differences between the fresh Nelder--Mead and
    stored solutions, and between the bounded TRF and fresh
    Nelder--Mead solutions, as a function of the stored mismatch. The
    sample spans the full mismatch distribution, including the
    $D_{\rm BSPL-PSPL}\simeq10^{-2}$ transition, parameter-bias extrema,
    and the boundaries of the sampled $(u_0,P)$ domain. Agreement at
    numerical precision demonstrates both reproducibility of the stored
    production minima and insensitivity to the optimizer when the
    objective function is held fixed.
    }
    \label{fig:stratified_refit_audit}
\end{figure*}

We also tested the temporal resolution by repeating representative
calculations with $5\times10^3$--$8\times10^4$ uniformly spaced epochs
over the same $t_0\pm3.5t_E$ integration window. For both a
single-luminous-source configuration with $u_0=10^{-2}$ and
$P\simeq t_E$, and a representative photocenter-cancellation
configuration with $q_f=q_M$, the production sampling of $10^4$ epochs
agrees with the $8\times10^4$-epoch result to better than $10^{-12}$ in
relative $D_{\rm BSPL-PSPL}$. This convergence applies to the broad
features used to interpret the parameter-space maps. Extremely close
point-source approaches can produce much narrower excursions and are
therefore treated with a separate local-refinement test in
Section~\ref{sec:close_approach}.



\subsection{Localized close-approach features}
\label{app:close_approach}

The point-source maps contain narrow structures associated with extreme
source--lens close approaches whose characteristic duration can be much
shorter than the global high-magnification timescale $u_0t_E$. In a
deliberately selected configuration with $q_M=0.09$ and $q_f=0.01$,
uniform grids containing $5\times10^3$--$8\times10^4$ epochs give
non-monotonic estimates of the mismatch ranging from
$D_{\rm BSPL-PSPL}=0.171$ to $0.602$. The production grid with
$10^4$ epochs therefore does not reliably resolve this localized
excursion.

A locally refined point-source calculation centered on the close
approach converges to $D_{\rm BSPL-PSPL}=0.569329$. This value is quoted
to six decimals here specifically to document the numerical convergence:
increasing both the number of locally sampled epochs and the width of
the refined temporal region changes the result by less than $10^{-6}$
fractionally. The discrepancy with the uniform grids is therefore caused
by temporal resolution rather than by convergence to a different PSPL
minimum.

We also performed targeted finite-source sanity checks designed to
separate genuine source-size effects from numerical under-resolution.
For the extreme close-approach configuration with $q_M=0.09$ and
$q_f=0.01$, the minimum lens--secondary-source separation is
$u_{2,\min}=1.71\times10^{-5}$. Over the sequence
$\rho_2=10^{-4},\,3\times10^{-4},\,10^{-3},\,3\times10^{-3}$, this
corresponds to $u_{2,\min}/\rho_2=0.171,\,0.0569,\,0.0171,$ and
$0.00569$, respectively. The finite-source calculation therefore probes
a regime in which the lens trajectory passes well within the projected
secondary-source radius. Consistently, the intrinsic mismatch decreases
monotonically from the converged point-source value
$D_{\rm BSPL-PSPL}=0.569329$ to $0.429$, $0.266$, $0.149$, and $0.086$
as $\rho_2$ is increased. For the largest radius considered,
$D_{\rm finite}/D_{\rm point}=0.151$, corresponding to a reduction by a
factor of about $6.6$. The peak residual is suppressed even more strongly,
while its characteristic duration increases, as expected when a very
narrow point-source excursion is smoothed and broadened by the finite
source.

As a control, we repeated the same sequence for $q_M=0.30$ and
$q_f=0.01$. In this case $u_{2,\min}=1.31\times10^{-2}$ and
$u_{2,\min}/\rho_2$ remains larger than $4$ throughout the explored
range. The mismatch is correspondingly unchanged: it varies only from
$0.0784629$ in the point-source calculation to $0.0784779$ at
$\rho_2=3\times10^{-3}$, i.e. by less than $2\times10^{-4}$
fractionally. The comparison therefore shows that the finite-source
response is localized to genuine close approaches rather than being a
generic numerical consequence of introducing non-zero source radii.

These tests also determine the qualitative direction in which finite-source
effects should modify the small-$u_0$ degeneracy boundary. When a sharp
point-source deviation is generated by $u_{2,\min}\lesssim\rho_2$,
finite-source smoothing reduces rather than enhances the intrinsic
BSPL--single-source mismatch. We therefore expect such effects to shift
configurations toward smaller $D_{\rm BSPL-PSPL}$ and hence to tend to
expand the region classified as degenerate near the
$D_{\rm BSPL-PSPL}=10^{-2}$ boundary. The present tests establish this
directional expectation but do not measure the displacement of the
boundary itself; a systematic remapping as a function of source radius
is beyond the scope of this work.

\subsection{Sensitivity to the integration window}
\label{app:window_sensitivity}

Because $D_{\rm BSPL-PSPL}$ is defined by a time integral over a finite
event-centered interval, its numerical value is conditional on the
adopted window. We therefore repeated representative intrinsic fits over
$t_0\pm2.5t_E$, $t_0\pm3.5t_E$, and $t_0\pm5t_E$, preserving
approximately the temporal sampling density of the production
calculation.

For the short- and intermediate-period configurations, the
high-magnification $u_0=10^{-2}$ case, and representative
two-luminous-source configurations, the fractional change in
$D_{\rm BSPL-PSPL}$ relative to the fiducial $3.5t_E$ window is below
$4\times10^{-5}$. The largest dependence occurs in the slowly varying
long-period cases. For the representative $P=6000\,{\rm d}$ case,
$D_{\rm BSPL-PSPL}$ changes by $-1.9\%$ for the $2.5t_E$ window and
$+0.6\%$ for the $5t_E$ window; the extreme hidden long-period case
changes by $-1.4\%$ and $+0.5\%$, respectively. In all three windows,
the representative intermediate-period mismatch remains larger than
both the short- and long-period values.

Figure~\ref{fig:window_sensitivity} summarizes these variations. The
test shows that the precise numerical value of the metric is defined
with respect to the adopted event window, while the broad physical
degeneracy structure is insensitive to reasonable changes in that
choice.

\begin{figure}[t]
    \centering
    \includegraphics[width=0.78\linewidth]{../figures/appendix/window_sensitivity.pdf}
    \caption{
    Sensitivity of the intrinsic mismatch to the event-centered
    integration window. The ordinate shows the fractional change in
    $D_{\rm BSPL-PSPL}$ relative to the fiducial
    $t_0\pm3.5t_E$ calculation. Broad short-, intermediate-, small-$u_0$,
    and two-luminous-source configurations change by less than
    $4\times10^{-5}$ fractionally. The largest response occurs for the
    slowly varying long-period cases, at the percent level, while the
    characteristic period ordering remains unchanged.
    }
    \label{fig:window_sensitivity}
\end{figure}
\subsection{Auxiliary solver-bound test}
\label{app:roman_te_bound}

The matched $t_E=30\,$d Roman map in the main text is not defined by the
following calculation. As an auxiliary solver-robustness test, we retain
the earlier $t_E=150\,$d grid-wide audit because it directly probes
whether truncating a formally weakly constrained PSPL timescale can
change the minimized Roman model-separation statistic. In that auxiliary
production we flagged configurations for which the fitted timescale lay
within $1\%$ of the upper bound,
\begin{equation}
t_{E,\mathrm{fit}}\geq0.99\,t_{E,\max},
\end{equation}
with $t_{E,\max}=20\,t_{E,\mathrm{true}}$. Only one physical
configuration showed this bound-limited behavior at all three source
magnitudes, at $u_0=0.977$ and $P/t_E=1.294$. This configuration lies
close to the upper edge of the retained Roman grid, $u_0\leq1$, but the
grid-wide diagnostic shows that the behavior is isolated within the
sampled domain: no other physical configuration is bound-limited at all
three magnitudes. We therefore repeated the fits for this configuration
while increasing the upper bound from $20\,t_{E,\rm true}$ to
$200\,t_{E,\rm true}$.

Figure~\ref{fig:roman_te_bound} shows that the fitted solution follows
an approximately flat PSPL direction in which $t_{E,\rm fit}$ increases
while $|u_{0,\rm fit}|$ decreases, leaving
$|u_{0,\rm fit}|t_{E,\rm fit}\simeq72.5$ d nearly unchanged. Despite
this parameter non-identifiability, the model-separation statistic is
already stable: between the $20\,t_{E,\rm true}$ and
$200\,t_{E,\rm true}$ bounds, $\Delta\chi^2_{\rm Roman}$ changes by only
$0.044\%$, $0.039\%$, and $0.030\%$ for W149$=19$, 21, and 23,
respectively. Thus, although the adopted production bound can truncate
the formal PSPL parameter solution in this configuration, its effect on
the survey-level model separation is negligible.

For the same auxiliary $t_E=150\,$d production, we applied the
$1\%$ proximity diagnostic to the lower bound,
$t_{E,\min}=0.02\,t_{E,\rm true}=3$ d. No fits lie
within $1\%$ of this lower bound at any of the three source magnitudes.
The smallest fitted timescales are $77.87$, $77.89$, and $77.94$ d for
W149$=19$, 21, and 23, respectively. Bound sensitivity is therefore
confined to the upper limit in the sampled Roman grid.

\begin{figure}[t]
    \centering
    \includegraphics[width=\linewidth]{../figures/appendix/roman_te_bound_validation.pdf}
    \caption{
    Sensitivity of the Roman PSPL fit to the adopted upper bound on
    $t_E$ for the configuration
    $u_0=0.977$ and $P/t_E=1.294$ that reaches the production bound.
    Panel (a) shows the best-fitting timescale as the upper bound is
    increased from $20$ to $200\,t_{E,\rm true}$; the dashed line marks
    equality with the imposed bound.
    Panel (b) shows the corresponding effective timescale
    $|u_{0,\rm fit}|t_{E,\rm fit}$, which remains nearly constant.
    Panel (c) shows the fractional change in
    $\Delta\chi^2_{\rm Roman}$ relative to the production
    $20\,t_{E,\rm true}$ fit.
    Although the formal PSPL parameters move along a nearly flat
    direction, the model-separation statistic changes by less than
    $5\times10^{-4}$ fractionally for all three source magnitudes.
    }
    \label{fig:roman_te_bound}
\end{figure}


\section{Robustness to projected orbital geometry}
\label{app:geometry_robustness}

The production scans in the main text adopt the projected orbital
geometry $(\alpha,\phi,i)=(0,0,\pi/2)$ in order to isolate the
dependence on the physical parameters being varied. The targeted geometry
validation below retains its original $t_E=150\,$d long-timescale setup;
it is used to test orientation dependence rather than to define the new
fiducial timescale. Since the detailed
BSPL light curve also depends on orbital orientation, we performed a
targeted robustness experiment to determine which conclusions are
specific to this fiducial projection and which persist under changes in
projected geometry.

We considered 60 projected orbital geometries. Of these, 48 form a
deterministic grid with
\begin{equation}
i\in\{0,\pi/4,\pi/2\},
\qquad
\phi\in\{0,\pi/2,\pi,3\pi/2\},
\end{equation}
and
\begin{equation}
\alpha\in\{0,\pi/4,\pi/2,3\pi/4\}.
\end{equation}
We supplemented this grid with 12 additional orientations sampled
randomly with uniform orbital phase and sky-plane orientation and
uniform $\cos i$. These calculations use the same
$t_0\pm3.5t_E$ integration window, $10^4$ temporal samples, and
re-optimized PSPL projection as the production scans.

\subsection{One-luminous-source period dependence}

We first repeated representative one-luminous-source calculations for
$q_M=0.5$, $u_0=0.1$, $t_E=150\,{\rm d}$, and
$\hat r_E=5\,{\rm AU}$ at $P=10$, $210$, and $6000\,{\rm d}$,
representing the short-, intermediate-, and long-period regimes.
Table~\ref{tab:geometry_broad} summarizes the distribution of
$D_{\rm BSPL-PSPL}$ across the 60 geometries. We also tested a
high-magnification configuration with $u_0=0.01$ and $P=210\,{\rm d}$.

\begin{table}[H]
\centering
\caption{
Dependence of the intrinsic mismatch on projected orbital geometry for
representative configurations. The fiducial column corresponds to the
geometry used in the production scans, while the remaining columns
summarize the distribution over all 60 tested orientations. The
two-luminous configurations use $q_M=0.5$, $P=t_E=150\,{\rm d}$, and
$q_f=0.5$ or $0.1$ as indicated.
}
\label{tab:geometry_broad}
\begin{tabular}{lcccc}
\toprule
Configuration &
$D_{\rm fid}$ &
$D_{16}$ &
$D_{\rm med}$ &
$D_{84}$ \\
\midrule
One luminous, $P=10\,{\rm d}$ &
$5.40\times10^{-2}$ &
$3.83\times10^{-2}$ &
$4.86\times10^{-2}$ &
$5.45\times10^{-2}$ \\
One luminous, $P=210\,{\rm d}$ &
$9.44\times10^{-2}$ &
$4.95\times10^{-2}$ &
$7.14\times10^{-2}$ &
$9.44\times10^{-2}$ \\
One luminous, $P=6000\,{\rm d}$ &
$2.11\times10^{-3}$ &
$9.74\times10^{-5}$ &
$1.26\times10^{-3}$ &
$1.86\times10^{-3}$ \\
One luminous, $u_0=0.01$, $P=210\,{\rm d}$ &
$3.42\times10^{-2}$ &
$1.79\times10^{-2}$ &
$2.49\times10^{-2}$ &
$3.54\times10^{-2}$ \\
Two luminous, $q_f=q_M=0.5$ &
$2.82\times10^{-1}$ &
$3.37\times10^{-2}$ &
$9.48\times10^{-2}$ &
$2.11\times10^{-1}$ \\
Two luminous, $q_f=0.1$ &
$1.37\times10^{-1}$ &
$5.23\times10^{-2}$ &
$8.98\times10^{-2}$ &
$1.10\times10^{-1}$ \\
\bottomrule
\end{tabular}
\end{table}

The absolute mismatch can vary substantially with orbital projection,
especially in the intermediate-period regime. Nevertheless, the
qualitative short--intermediate--long-period structure remains stable:
in 54 of the 60 tested validation configurations, the
intermediate-period configuration has a larger mismatch than both the
short- and long-period configurations. Because this validation set
combines a deterministic angular grid with a smaller isotropically
sampled subset, the fraction $54/60$ is a validation-set diagnostic and
should not be interpreted as an orientation-weighted probability. The
exceptions correspond to orientations for which the projected orbital
perturbation is additionally suppressed. More importantly, the long-period configuration remains
strongly degenerate for every tested orientation; across the full
sample its mismatch never exceeds approximately
$2.2\times10^{-3}$. Thus the detailed contours in the fiducial maps are
geometry dependent, whereas the existence of a strongly degenerate
long-period regime is not. This behavior is shown directly in
Figure~\ref{fig:geometry_robustness}A, where the three representative
period regimes are connected for each tested orbital geometry.

\subsection{Orientation independence of the photocenter scaling}

The photocenter-cancellation argument makes a stronger prediction than
the finite-separation maps: when $q_f=q_M$, the first-order term in the
binary-source deformation vanishes identically, independent of the
direction of the projected displacement. We tested this explicitly by
repeating the small-separation experiment of
Section~\ref{sec:photocenter_cancellation} for all 60 geometries.

For each orientation we fixed
$q_M=0.5$, $u_0=0.1$, and $P=t_E=150\,{\rm d}$. We then varied
$\xi_{\rm rel}/u_0$ over nine logarithmically spaced values between
$10^{-4}$ and $10^{-2}$, independently of the Keplerian
period--separation relation, by rescaling $\hat r_E$ while keeping the
orbital time dependence fixed. We fitted
\begin{equation}
D_{\rm BSPL-PSPL}\propto\xi_{\rm rel}^{\gamma}.
\end{equation}
The resulting logarithmic slopes are summarized in
Table~\ref{tab:geometry_slopes}.

\begin{table}[H]
\centering
\caption{
Logarithmic slope $\gamma$ of the small-separation relation
$D_{\rm BSPL-PSPL}\propto\xi_{\rm rel}^{\gamma}$ across the 60 tested
orbital geometries. The quoted interval gives the 16th--84th percentile
range, and the final column gives the full range.
}
\label{tab:geometry_slopes}
\begin{tabular}{lccc}
\toprule
Source configuration &
Median $\gamma$ &
16th--84th percentile &
Full range \\
\midrule
Dark companion ($q_f=0$) &
$1.000000$ &
$0.999868$--$1.000132$ &
$0.999804$--$1.000196$ \\
Photocenter cancellation ($q_f=q_M$) &
$2.000001$ &
$1.998876$--$2.001094$ &
$1.998137$--$2.001874$ \\
\bottomrule
\end{tabular}
\end{table}

All fits have $R^2$ numerically indistinguishable from unity, and every
tested geometry lies within $0.1$ of the expected slope. The
orientation dependence therefore changes the amplitude of the
perturbation but not its leading asymptotic order:
\begin{equation}
D_{\rm BSPL-PSPL}=\mathcal{O}(\xi_{\rm rel})
\qquad (q_f=0),
\end{equation}
whereas
\begin{equation}
D_{\rm BSPL-PSPL}=\mathcal{O}(\xi_{\rm rel}^2)
\qquad (q_f=q_M).
\end{equation}
This provides a direct numerical check that the first-order
photocenter cancellation is not tied to the fiducial orbital
orientation. Figure~\ref{fig:geometry_robustness}B shows the deviations
of the fitted slopes from their analytic expectations,
$\gamma_{\rm dark}=1$ and $\gamma_{\rm cancel}=2$, for all 60
orientations.

At finite separation the cancellation condition should not be
interpreted as requiring
$D(q_f=q_M)<D(q_f\neq q_M)$ for every geometry. In the representative
$P=t_E$ comparison, the median ratio
$D(q_f=q_M)/D(q_f=0.1)$ is $0.99$, with the cancellation case having
the smaller mismatch in 31 of 60 orientations. This is expected because
the analytic cancellation concerns the leading term in the
small-$\xi_{\rm rel}$ expansion; when $\xi_{\rm rel}/u_0$ is of order
unity, quadratic and higher-order terms can dominate the total
mismatch. The robust result is therefore the removal of the linear term,
not a universal ordering of finite-separation values of
$D_{\rm BSPL-PSPL}$.



\begin{figure*}[t]
    \centering
    \includegraphics[width=\textwidth]{../figures/appendix/geometry_robustness_appendix.pdf}
\caption{
Robustness of the intrinsic BSPL--PSPL mismatch to projected orbital
geometry.
\textit{Left:} Distribution of $D_{\rm BSPL-PSPL}$ for representative
short-, intermediate-, and long-period one-luminous-source
configurations across the 60 orbital geometries considered in
Appendix~\ref{app:geometry_robustness}. Open points show individual
orientations, red symbols and error bars show the median and
16th--84th percentile interval, and black stars indicate the fiducial
production geometry. The intermediate-period mismatch exceeds both the
short- and long-period values in 54 of the 60 tested validation
configurations, while the long-period configuration remains strongly
degenerate in every tested case; $54/60$ is a validation-set fraction,
not an orientation-weighted probability.
\textit{Right:} Deviations of the fitted small-separation power-law
slopes from their analytic expectations. Individual points correspond
to the 60 orbital orientations, while diamonds and error bars show the
median and 16th--84th percentile interval. The dark-companion family
remains consistent with $\gamma=1$, whereas the photocenter-cancellation
family remains consistent with $\gamma=2$ for all tested orientations,
demonstrating that the change from linear to quadratic asymptotic
scaling is independent of projected orbital geometry.
}
    \label{fig:geometry_robustness}
\end{figure*}


\section*{Data and code availability}

The analysis code is maintained in the public repository
\url{https://github.com/anibal-art/ulensing_degenerate_models}. The
repository contains the intrinsic and Roman production scripts,
validation workflows, machine-readable summary products, and figure
producers used in this work. The principal numerical products store code
provenance together with the fitted grids. The PARSEC/CMD isochrone table
used to construct the Roman F146 mass--flux tracks is retained with the
stellar-physics analysis products. A tagged archival release containing
the exact manuscript source, derived data products, and software-version
information will be created at submission. The numerical setup needed to
reproduce the principal intrinsic and Roman calculations is summarized in
Tables~\ref{tab:numerical_setup} and~\ref{tab:roman_setup}.

\bibliographystyle{plainnat}
\bibliography{latex/sample}

\end{document}
